\documentclass[aspectratio=169]{beamer}

% =========================================================
% Кодировка и язык
% =========================================================
\usepackage[utf8]{inputenc}
\usepackage[T2A]{fontenc}
\usepackage[russian]{babel}

% =========================================================
% Математика
% =========================================================
\usepackage{amsmath,amssymb}
\usepackage{bm}

% =========================================================
% Графика и таблицы
% =========================================================
\usepackage{graphicx}
\usepackage{booktabs}
\usepackage{array}

% =========================================================
% TikZ (схемы, архитектуры, диаграммы)
% =========================================================
\usepackage{tikz}
\usetikzlibrary{shapes.geometric, arrows.meta, positioning, calc, fit, backgrounds}

% =========================================================
% Тема
% =========================================================
\usetheme{default}
\usecolortheme{default}

% =========================================================
% Цвета СВФУ
% =========================================================
\definecolor{svfuDark}{RGB}{10,40,90}
\definecolor{svfuAccent}{RGB}{30,144,255}
\definecolor{svfuLight}{RGB}{230,240,255}

% =========================================================
% Стили TikZ
% =========================================================
\tikzset{
  entity/.style={
    rectangle, draw=svfuDark, very thick, fill=svfuLight,
    minimum width=2.4cm, minimum height=1cm,
    font=\bfseries\small, align=center, rounded corners=1pt
  },
  relationship/.style={
    diamond, draw=svfuDark, very thick, fill=svfuAccent!25,
    aspect=2.4, font=\scriptsize, align=center, inner sep=1pt
  },
  attribute/.style={
    ellipse, draw=svfuDark, thin, fill=white,
    minimum width=1.8cm, minimum height=0.55cm,
    font=\scriptsize, align=center
  },
  card/.style={font=\tiny, fill=svfuLight, inner sep=1pt},
  layer/.style={
    rectangle, draw=svfuDark, very thick, fill=svfuAccent!15,
    minimum width=6.2cm, minimum height=1.3cm,
    font=\bfseries\normalsize, align=center, rounded corners=2pt
  },
  phase/.style={
    rectangle, draw=svfuDark, thick, fill=svfuLight,
    rounded corners=4pt, minimum width=2.5cm, minimum height=1cm,
    font=\small\bfseries, align=center
  },
  catbox/.style={
    rectangle, draw=svfuDark, thick, fill=white,
    minimum width=2.6cm, minimum height=0.9cm,
    font=\footnotesize, align=center, rounded corners=1pt
  },
  module/.style={
    rectangle, draw=svfuDark, very thick, fill=white,
    minimum width=3.4cm, minimum height=1.1cm,
    font=\footnotesize\bfseries, align=center, rounded corners=2pt
  },
  actor/.style={
    rectangle, draw=svfuDark, thick, fill=svfuLight,
    minimum width=2.6cm, minimum height=1cm,
    font=\small\bfseries, align=center, rounded corners=8pt
  }
}

% =========================================================
% Основные цвета Beamer
% =========================================================
\setbeamercolor{normal text}{fg=black,bg=white}
\setbeamercolor{structure}{fg=svfuDark}
\setbeamercolor{alerted text}{fg=svfuDark}
\setbeamercolor{frametitle}{bg=svfuDark,fg=white}
\setbeamercolor{title}{fg=white}
\setbeamercolor{author}{fg=white}
\setbeamercolor{date}{fg=white}

\setbeamercolor{accent line}{bg=svfuAccent}

\setbeamercolor{footline left}
  {bg=svfuDark,fg=white}

\setbeamercolor{footline center}
  {bg=svfuDark!85,fg=white!80}

\setbeamercolor{footline right}
  {bg=svfuDark!70,fg=white}

\setbeamercolor{block title}
  {bg=svfuDark,fg=white}

\setbeamercolor{block body}
  {bg=svfuLight,fg=black}

% =========================================================
% Шрифты
% =========================================================
\setbeamerfont{frametitle}
  {size=\large,series=\bfseries}

\setbeamerfont{title}
  {size=\Large,series=\bfseries}

% =========================================================
% Убираем навигацию
% =========================================================
\setbeamertemplate{navigation symbols}{}
\setbeamertemplate{headline}{}

% =========================================================
% Заголовок слайда
% =========================================================
\setbeamertemplate{frametitle}{%
  \vspace*{-1pt}%
  \begin{beamercolorbox}[
    wd=\paperwidth,
    ht=3.2ex,
    dp=1.2ex,
    leftskip=10pt,
    sep=0pt
  ]{frametitle}
    \usebeamerfont{frametitle}
    \insertframetitle
  \end{beamercolorbox}%
}

% =========================================================
% Подвал
% =========================================================
\setbeamertemplate{footline}{%
  \begin{beamercolorbox}[
    wd=\paperwidth,
    ht=0.4pt,
    dp=0pt
  ]{accent line}
  \end{beamercolorbox}%

  \hbox to \paperwidth{%

    \begin{beamercolorbox}[
      wd=0.305\paperwidth,
      ht=3ex,
      dp=1.5ex,
      left,
      sep=0pt,
      leftskip=6ex
    ]{footline left}
      \textbf{СВФУ}
    \end{beamercolorbox}%

    \begin{beamercolorbox}[
      wd=0.295\paperwidth,
      ht=3ex,
      dp=1.5ex,
      center,
      sep=0pt
    ]{footline center}
    \end{beamercolorbox}%

    \begin{beamercolorbox}[
      wd=0.4\paperwidth,
      ht=3ex,
      dp=1.5ex,
      right,
      sep=0pt
    ]{footline right}
      \raggedleft
      \today
      \hspace{5ex}
      \insertframenumber\,/\,\inserttotalframenumber
      \hspace{6ex}
    \end{beamercolorbox}%
  }%
}

% =========================================================
% Настройки списков
% =========================================================
\setbeamertemplate{itemize item}
  {\color{svfuAccent}\small$\blacktriangleright$}

\setbeamertemplate{itemize subitem}
  {\color{svfuAccent}\scriptsize$\blacktriangleright$}

\setlength{\leftmargini}{1.5em}

% =========================================================
% Документ
% =========================================================
\begin{document}

% =========================================================
% 1. ТИТУЛ
% =========================================================
\begin{frame}[plain]
  \vspace*{1cm}
  \begin{center}

    \colorbox{svfuDark}{%
      \parbox[c][2.8cm][c]{0.86\paperwidth}{%
        \centering
        {\LARGE\color{white}
          Лекция 2. Системы управления базами данных (СУБД)
        }
      }%
    }

  \end{center}

  \vfill

  \centering
  {\normalsize Дисциплина: Методы и средства проектирования баз данных}\\[4pt]
  {\small
    Составитель: Монастырев Егорий Егорьевич, М-ВТ-25
  }\\[4pt]

  {\small
    Якутск, 2026
  }

  \vfill

\end{frame}


% =========================================================
% СВЯЗЬ С ЛЕКЦИЕЙ 1 + ЦЕЛИ (объединённый слайд)
% =========================================================
\begin{frame}{Связь с предыдущей лекцией и цели сегодняшней}

\small
\textbf{Уже знаем (Лекция 1):} данные и модели их представления,
предметная область, ER-модель (сущность, атрибут, связь).

\vspace{0.4cm}

\textbf{Сегодня разбираемся:}
\begin{enumerate}
  \item Чем база данных, СУБД и система базы данных отличаются друг от друга?
  \item Какие функции обязана выполнять СУБД и зачем?
  \item Из каких программных компонентов состоит СУБД изнутри?
  \item Как СУБД обслуживает многих пользователей одновременно?
  \item Из каких подъязыков состоит SQL и как ими пользоваться на практике?
\end{enumerate}

\vspace{0.3cm}

\begin{alertblock}{Зачем это нужно уже сейчас}
На ближайшей лабораторной работе вы начнёте писать
реальные SQL-команды — сегодняшняя лекция даёт для этого
необходимый минимум.
\end{alertblock}

\end{frame}


% =========================================================
% РАЗДЕЛ 1. ОТ БАЗЫ ДАННЫХ К СИСТЕМЕ УПРАВЛЕНИЯ БАЗАМИ ДАННЫХ
% =========================================================

\begin{frame}{Почему одной базы данных недостаточно}

Представим, что данные просто лежат в файлах на диске,
без управляющей программы.

\vspace{0.3cm}

\begin{itemize}
  \item Кто гарантирует, что два сотрудника не испортят
        одну и ту же запись, редактируя её одновременно?
  \item Кто восстановит данные после сбоя электропитания
        посреди записи?
  \item Кто проверит, что оценка не выставлена
        несуществующему студенту?
  \item Кто ограничит доступ так, чтобы студент не мог
        править чужие оценки?
\end{itemize}

\vspace{0.3cm}

\begin{alertblock}{Вывод}
Сами данные не решают эти задачи. Ими управляет
специализированное программное обеспечение — \textbf{СУБД}.
\end{alertblock}

\end{frame}


\begin{frame}{База данных и СУБД: уточнённые определения}

\small
\begin{columns}[T]
\column{0.5\textwidth}
\begin{block}{База данных}
Организованная по \textbf{схеме} совокупность
взаимосвязанных данных, находящаяся под управлением
специализированного ПО, предназначенная для совместного
использования многими приложениями и пользователями.
\end{block}
\column{0.5\textwidth}
\begin{block}{СУБД}
Комплекс программных средств, реализующий определение
схемы, обработку транзакций, поддержание целостности,
разграничение доступа, журналирование и восстановление.
\end{block}
\end{columns}

\vspace{0.4cm}

\begin{alertblock}{Почему не просто «набор данных» / «программа для БД»}
Такие бытовые формулировки не отличают базу данных
от таблицы Excel и не объясняют, зачем именно нужна СУБД —
целостность, многопользовательский доступ, восстановление.
\end{alertblock}

\end{frame}


\begin{frame}{От СУБД к системе базы данных}

\begin{columns}[T]

\column{0.52\textwidth}
\centering
\resizebox{\linewidth}{!}{%
\begin{tikzpicture}[node distance=0.3cm and 0.5cm]

\node[entity, minimum width=4.6cm, minimum height=0.75cm] (data) {База данных};
\node[layer, minimum width=4.6cm, minimum height=0.75cm, below=of data] (dbms) {СУБД};
\node[actor, fill=white, minimum height=0.7cm, below=of dbms, xshift=-1.3cm] (apps) {Приложения};
\node[actor, fill=white, minimum height=0.7cm, below=of dbms, xshift=1.3cm] (people) {Пользователи};

\begin{scope}[on background layer]
\node[draw=svfuAccent, very thick, rounded corners=6pt,
      fit=(data)(dbms)(apps)(people), inner sep=0.28cm,
      label={[font=\bfseries\scriptsize]above:Система базы данных}]
      (frame) {};
\end{scope}

\end{tikzpicture}%
}

\column{0.44\textwidth}
\scriptsize
\begin{tabular}{>{\bfseries}p{0.34\textwidth}|p{0.55\textwidth}}
\toprule
Понятие & Что это\\
\midrule
База данных & данные + схема\\
СУБД & программа управления данными\\
Система БД & данные + СУБД + приложения + люди\\
\bottomrule
\end{tabular}

\end{columns}

\vspace{0.35cm}
\small
\begin{alertblock}{На примере «Университета»}
Данные о студентах — \textbf{БД}; PostgreSQL — \textbf{СУБД};
портал деканата (БД + СУБД + приложение + сотрудники) —
\textbf{система базы данных}.
\end{alertblock}

\end{frame}


% =========================================================
% РАЗДЕЛ 2. НАЗНАЧЕНИЕ И ФУНКЦИИ СУБД
% =========================================================

\begin{frame}{Общее назначение СУБД}

СУБД решает задачи, которые невозможно надёжно решить
средствами отдельных файлов и прикладных программ:

\vspace{0.3cm}

\begin{itemize}
  \item предоставляет единый, контролируемый способ
        доступа к данным;
  \item берёт на себя ответственность за корректность
        данных при сбоях и параллельной работе;
  \item скрывает от приложений детали физического
        хранения данных.
\end{itemize}

\vspace{0.3cm}

\begin{block}{Идея}
Приложения работают с данными \textbf{через} СУБД,
а не напрямую с файлами на диске.
\end{block}

\end{frame}


\begin{frame}{Функции СУБД}

\scriptsize
\begin{table}
\centering
\begin{tabular}{p{0.24\textwidth}|p{0.34\textwidth}|p{0.32\textwidth}}
\toprule
\textbf{Функция} & \textbf{Суть} & \textbf{Пример («Университет»)}\\
\midrule
Хранение & таблицы, файлы, страницы & данные о студентах на диске\\
Обработка запросов & разбор и оптимизация запросов & поиск студента по фамилии\\
Транзакции & «всё или ничего» & перевод студента в другую группу\\
Целостность & проверка правил корректности & оценка не привязана к несуществующему студенту\\
Многопольз. доступ & корректная одновременная работа & два сотрудника правят одну запись\\
Журналирование и восстановление & защита от потери данных & сбой питания не стирает ведомость\\
Безопасность & разграничение прав доступа & студент видит только свои оценки\\
\bottomrule
\end{tabular}
\end{table}

\end{frame}


\begin{frame}{От функций СУБД — к свойствам базы данных}

\begin{center}
\begin{tikzpicture}[node distance=0.22cm and 1.2cm]

\node[module, fill=svfuLight, minimum height=0.68cm] (f1) {Целостность};
\node[module, fill=svfuLight, minimum height=0.68cm, below=of f1] (f2) {Многопольз.\\доступ};
\node[module, fill=svfuLight, minimum height=0.68cm, below=of f2] (f3) {Безопасность};
\node[module, fill=svfuLight, minimum height=0.68cm, below=of f3] (f4) {Хранение};

\node[catbox, minimum height=0.68cm, right=2.4cm of f1] (p1) {Целостность (БД)};
\node[catbox, minimum height=0.68cm, right=2.4cm of f2] (p2) {Разделяемость};
\node[catbox, minimum height=0.68cm, right=2.4cm of f3] (p3) {Защищённость};
\node[catbox, minimum height=0.68cm, right=2.4cm of f4] (p4) {Интегрированность};

\draw[-{Latex[length=2mm]}, thick, svfuAccent] (f1) -- (p1);
\draw[-{Latex[length=2mm]}, thick, svfuAccent] (f2) -- (p2);
\draw[-{Latex[length=2mm]}, thick, svfuAccent] (f3) -- (p3);
\draw[-{Latex[length=2mm]}, thick, svfuAccent] (f4) -- (p4);

\node[above=0.15cm of f1, font=\footnotesize\bfseries] {Функция СУБД};
\node[above=0.15cm of p1, font=\footnotesize\bfseries] {Свойство БД};

\end{tikzpicture}
\end{center}

\small
\begin{alertblock}{Главная мысль}
Свойства базы данных, названные в Лекции 1
(целостность, разделяемость, защищённость,
интегрированность), — не сами по себе, а \textbf{результат}
работы конкретных функций СУБД.
\end{alertblock}
\normalsize

\end{frame}


% =========================================================
% РАЗДЕЛ 3. АРХИТЕКТУРА ПРЕДСТАВЛЕНИЯ ДАННЫХ
% =========================================================

\begin{frame}{Схема и экземпляр базы данных}

\begin{itemize}
  \item \textbf{Схема базы данных} — формальное описание
        структуры БД: таблицы, атрибуты, типы данных,
        связи, ограничения. Меняется редко.
  \item \textbf{Экземпляр базы данных} — конкретное состояние
        данных в определённый момент времени. Меняется
        практически при каждой операции записи.
\end{itemize}

\vspace{0.4cm}

\begin{alertblock}{Аналогия для запоминания}
Схема — это «тип» (структура), экземпляр — это «значение»
(конкретное содержимое) в данный момент.
\end{alertblock}

\vspace{0.2cm}
\footnotesize
Пример: схема таблицы \texttt{student} (id, ФИО, group\_id)
не меняется годами; строки этой таблицы меняются
каждый день.

\end{frame}


\begin{frame}{Архитектура ANSI/SPARC и независимость данных}

\footnotesize
Кратко напомним схему из Лекции 1 и свяжем её
с независимостью данных.

\begin{columns}[T]

\column{0.5\textwidth}
\centering
\resizebox{\linewidth}{!}{%
\begin{tikzpicture}[node distance=0.35cm]
\node[layer, minimum width=4.4cm, minimum height=1cm] (ext) {Внешний\\\footnotesize представления пользователей};
\node[layer, minimum width=4.4cm, minimum height=1cm, below=of ext] (con) {Концептуальный\\\footnotesize структура БД};
\node[layer, minimum width=4.4cm, minimum height=1cm, below=of con] (int) {Внутренний\\\footnotesize физическое хранение};
\draw[-{Latex[length=2.5mm]}, thick, svfuAccent] (ext) -- (con);
\draw[-{Latex[length=2.5mm]}, thick, svfuAccent] (con) -- (int);
\end{tikzpicture}%
}

\column{0.47\textwidth}
\scriptsize
\begin{tabular}{p{0.28\textwidth}|p{0.6\textwidth}}
\toprule
Тип независимости & Что позволяет\\
\midrule
Физическая & менять хранение и индексы без переписывания приложений\\
Логическая & менять концептуальную схему, сохраняя внешние представления\\
\bottomrule
\end{tabular}

\vspace{0.3cm}
Цель всей архитектуры — \textbf{независимость данных}.

\end{columns}

\vspace{0.25cm}
\begin{alertblock}{Важная оговорка}
Логическая независимость имеет пределы: удаление
атрибута, уже использованного приложением, всё равно
нарушит его работу.
\end{alertblock}

\end{frame}


% =========================================================
% РАЗДЕЛ 4. ВНУТРЕННЯЯ АРХИТЕКТУРА СУБД
% =========================================================

\begin{frame}{ANSI/SPARC и внутренняя архитектура — не одно и то же}

\begin{block}{Важное разграничение}
Архитектура ANSI/SPARC отвечает на вопрос
\textbf{«как данные видны снаружи»}.\\[4pt]
Внутренняя архитектура СУБД отвечает на вопрос
\textbf{«из каких программных модулей СУБД состоит изнутри»}.
\end{block}

\vspace{0.4cm}

Это разные по природе понятия, которые не стоит путать:
первое — про уровни абстракции представления данных,
второе — про устройство программного ядра.

\end{frame}


\begin{frame}{Компоненты ядра СУБД}

\small
\begin{table}
\centering
\begin{tabular}{p{0.42\textwidth}|p{0.44\textwidth}}
\toprule
\textbf{Компонент} & \textbf{Что делает}\\
\midrule
Обработчик запросов & разбор и оптимизация запросов\\
Менеджер транзакций & контроль атомарности и изоляции операций\\
Менеджер хранения и буфера & управление файлами, страницами, кэшем в ОЗУ\\
Менеджер журнала и восстановления & фиксация изменений, восстановление после сбоев\\
Словарь данных / системный каталог & хранение метаданных о структуре БД\\
\bottomrule
\end{tabular}
\end{table}

\vspace{0.3cm}
\footnotesize
Каждый компонент отвечает за одну из функций СУБД,
рассмотренных в разделе о назначении СУБД.

\end{frame}


\begin{frame}{Путь одного запроса через компоненты СУБД}

\begin{center}
\resizebox{0.98\textwidth}{!}{%
\begin{tikzpicture}[node distance=0.55cm and 0.7cm]

\node[actor, fill=svfuLight] (app) {Приложение\\\texttt{SELECT ...}};
\node[module, right=of app] (parser) {Обработчик\\запросов};
\node[module, right=of parser] (tx) {Менеджер\\транзакций};
\node[module, right=of tx] (storage) {Менеджер\\хранения};
\node[entity, right=of storage, minimum height=1.1cm] (disk) {Данные\\на диске};

\draw[-{Latex[length=2mm]}, thick, svfuAccent] (app) -- (parser);
\draw[-{Latex[length=2mm]}, thick, svfuAccent] (parser) -- (tx);
\draw[-{Latex[length=2mm]}, thick, svfuAccent] (tx) -- (storage);
\draw[-{Latex[length=2mm]}, thick, svfuAccent] (storage) -- (disk);

\node[module, below=1.1cm of tx] (log) {Менеджер\\журнала};
\draw[-{Latex[length=2mm]}, thick, svfuDark, dashed] (tx) -- (log);

\node[catbox, below=1.1cm of parser] (dict) {Словарь\\данных};
\draw[-{Latex[length=2mm]}, thick, svfuDark, dashed] (parser) -- (dict);

\end{tikzpicture}%
}
\end{center}

\footnotesize
Схема иллюстрирует общий путь запроса без разбора
алгоритмов оптимизации: обработчик обращается к словарю
данных, чтобы узнать структуру таблиц; менеджер транзакций
взаимодействует с менеджером журнала для гарантии
устойчивости изменений.

\end{frame}


% =========================================================
% РАЗДЕЛ 5. МНОГОПОЛЬЗОВАТЕЛЬСКИЙ РЕЖИМ
% =========================================================

\begin{frame}{Клиент-серверная модель доступа к СУБД}

\begin{center}
\begin{tikzpicture}[node distance=0.4cm and 0.9cm]

\node[actor, fill=white, minimum width=2.4cm] (c1) {Деканат};
\node[actor, fill=white, minimum width=2.4cm, below=0.35cm of c1] (c2) {Приёмная комиссия};
\node[actor, fill=white, minimum width=2.4cm, below=0.35cm of c2] (c3) {Личный кабинет};

\node[layer, right=1.8cm of c2, minimum width=3.6cm] (server) {Сервер СУБД\\\footnotesize (PostgreSQL)};

\node[entity, right=1.8cm of server, minimum height=1.4cm, minimum width=2.2cm] (files) {Файлы\\базы данных};

\draw[thick, svfuAccent] (c1) -- (server);
\draw[thick, svfuAccent] (c2) -- (server);
\draw[thick, svfuAccent] (c3) -- (server);
\draw[-{Latex[length=2.5mm]}, thick, svfuDark] (server) -- (files);

\end{tikzpicture}
\end{center}

\begin{alertblock}{Идея}
Множество клиентских приложений одновременно подключаются
к \textbf{одному} серверу СУБД, а не к файлам напрямую.
\end{alertblock}

\end{frame}


\begin{frame}{Параллельная работа: постановка проблемы}

\begin{center}
\large
Сотрудник А редактирует оценку студента 101
\end{center}
\begin{center}
\large
Сотрудник Б \textbf{одновременно} редактирует ту же оценку
\end{center}

\vspace{0.4cm}

\begin{itemize}
  \item Что произойдёт, если оба сохранят изменения
        одновременно?
  \item Как СУБД не допускает потери одного из изменений?
\end{itemize}

\vspace{0.3cm}

\begin{alertblock}{На этом этапе — только постановка вопроса}
Механика решения (блокировки, уровни изоляции)
рассматривается в отдельной будущей лекции
о транзакциях и параллельном выполнении. Сейчас важно
зафиксировать: СУБД обязана это контролировать.
\end{alertblock}

\end{frame}


\begin{frame}{Роли участников работы с СУБД}

\small
\begin{table}
\centering
\begin{tabular}{>{\bfseries}p{0.32\textwidth}|p{0.55\textwidth}}
\toprule
Роль & Что делает\\
\midrule
Администратор БД (АБД) &
администрирует безопасность, резервное копирование
и восстановление, следит за производительностью\\
Прикладные программисты &
используют транзакции и целостность через приложение,
не работая с ядром СУБД напрямую\\
Конечные пользователи &
взаимодействуют с данными через приложение,
не видя внутреннего устройства СУБД\\
\bottomrule
\end{tabular}
\end{table}

\end{frame}


% =========================================================
% РАЗДЕЛ 6. СУБД ДЛЯ ЛАБОРАТОРНЫХ РАБОТ
% =========================================================

\begin{frame}{СУБД для лабораторных работ: PostgreSQL и инструменты}

В курсе используется \textbf{PostgreSQL} — серверная СУБД
(в отличие от настольных СУБД вроде Access, уже
упомянутых в Лекции 1: ограниченный многопользовательский
доступ, локальное использование).

\vspace{0.4cm}

Подключение и работа с сервером PostgreSQL — через
отдельные клиентские инструменты:

\begin{itemize}
  \item \textbf{pgAdmin} — графический интерфейс: таблицы,
        запросы, права доступа;
  \item \textbf{psql} — консольный (CLI) клиент для прямой
        работы с сервером.
\end{itemize}

\vspace{0.4cm}

\begin{alertblock}{На лабораторных работах}
Именно через psql / pgAdmin вы будете выполнять
SQL-команды, которые рассматриваются далее в этой лекции.
\end{alertblock}

\end{frame}


% =========================================================
% РАЗДЕЛ 7. SQL КАК СРЕДСТВО РАБОТЫ С СУБД
% =========================================================

\begin{frame}{SQL как декларативный язык}

\begin{block}{Определение}
\textbf{SQL} (Structured Query Language) — стандартизированный
декларативный язык для определения структуры данных,
выполнения операций над данными и управления доступом
в реляционных СУБД.
\end{block}

\vspace{0.3cm}

\begin{itemize}
  \item декларативный — пользователь описывает,
        \textbf{что} нужно получить, а не \textbf{как} именно
        это будет вычислено;
  \item поддерживается практически всеми реляционными СУБД
        (с диалектными отличиями; далее — синтаксис PostgreSQL).
\end{itemize}

\end{frame}


\begin{frame}{Подъязыки SQL: карта языка}

\small
\begin{table}
\centering
\begin{tabular}{>{\bfseries}p{0.13\textwidth}|p{0.37\textwidth}|p{0.36\textwidth}}
\toprule
Подъязык & Назначение & Связь с функцией СУБД\\
\midrule
DDL & определение структуры данных & хранение\\
DML & работа с данными & обработка запросов\\
DCL & управление доступом & безопасность\\
TCL & управление транзакциями & транзакции\\
\bottomrule
\end{tabular}
\end{table}

\vspace{0.4cm}

\begin{alertblock}{Главная идея}
SQL — не хаотичный набор команд, а язык, структурно
повторяющий функции самой СУБД. Далее — каждый подъязык
подробнее, с рабочим синтаксисом для лабораторной работы.
\end{alertblock}

\end{frame}


\begin{frame}{DDL — язык определения данных}

\textbf{DDL (Data Definition Language)} — команды, которые
определяют и изменяют \textbf{структуру} базы данных:
таблицы, столбцы, ограничения.

\vspace{0.3cm}

\small
\begin{table}
\centering
\begin{tabular}{>{\ttfamily}p{0.22\textwidth}|p{0.62\textwidth}}
\toprule
\textnormal{Команда} & \textnormal{Назначение}\\
\midrule
CREATE & создать объект (таблицу, индекс и т.\,д.)\\
ALTER & изменить структуру существующего объекта\\
DROP & удалить объект\\
\bottomrule
\end{tabular}
\end{table}

\vspace{0.3cm}
\normalsize
\begin{alertblock}{Связь с функциями СУБД}
DDL — практический инструмент функции \textbf{«хранение
данных»}: этими командами задаётся схема, которую
использует СУБД.
\end{alertblock}

\end{frame}


\begin{frame}[fragile]{DDL — пример: создание таблиц}

\small
\begin{verbatim}
CREATE TABLE student_group (
    id      SERIAL PRIMARY KEY,
    name    TEXT NOT NULL
);

CREATE TABLE student (
    id         SERIAL PRIMARY KEY,
    full_name  TEXT NOT NULL,
    group_id   INTEGER REFERENCES student_group(id)
);
\end{verbatim}
\normalsize

\vspace{0.2cm}
\footnotesize
Ограничения (\texttt{PRIMARY KEY}, \texttt{REFERENCES})
подробно рассматриваются в Лекции 3; здесь важно увидеть,
как DDL превращает логическую структуру в реальные таблицы.

\end{frame}


\begin{frame}{DML — язык управления данными}

\textbf{DML (Data Manipulation Language)} — команды для
чтения и изменения \textbf{содержимого} таблиц.

\vspace{0.3cm}

\small
\begin{table}
\centering
\begin{tabular}{>{\ttfamily}p{0.22\textwidth}|p{0.62\textwidth}}
\toprule
\textnormal{Команда} & \textnormal{Назначение}\\
\midrule
SELECT & выбрать данные, удовлетворяющие условию\\
INSERT & добавить новую строку\\
UPDATE & изменить значения в существующих строках\\
DELETE & удалить строки\\
\bottomrule
\end{tabular}
\end{table}

\vspace{0.3cm}
\normalsize
\begin{alertblock}{Связь с функциями СУБД}
DML напрямую задействует \textbf{обработчик запросов}
(раздел о внутренней архитектуре) и опирается на
\textbf{целостность}, которую контролирует СУБД.
\end{alertblock}

\end{frame}


\begin{frame}[fragile]{DML — примеры базовых операций}

\small
\begin{verbatim}
SELECT full_name FROM student
WHERE group_id = 1;

INSERT INTO student (full_name, group_id)
VALUES ('Иванов И. И.', 1);

UPDATE student SET group_id = 2
WHERE id = 5;

DELETE FROM student WHERE id = 7;
\end{verbatim}
\normalsize

\vspace{0.2cm}
\footnotesize
Это базовый синтаксис для одной таблицы — достаточный
минимум для первой лабораторной работы.

\end{frame}


\begin{frame}[fragile]{DCL — язык управления доступом}

\textbf{DCL (Data Control Language)} — команды для
разграничения прав доступа к данным.

\vspace{0.3cm}

\small
\begin{table}
\centering
\begin{tabular}{>{\ttfamily}p{0.22\textwidth}|p{0.62\textwidth}}
\toprule
\textnormal{Команда} & \textnormal{Назначение}\\
\midrule
GRANT & предоставить право на объект (роли, пользователю)\\
REVOKE & отозвать ранее предоставленное право\\
\bottomrule
\end{tabular}
\end{table}

\vspace{0.3cm}

\begin{verbatim}
GRANT SELECT ON student TO dean_office;
REVOKE UPDATE ON student FROM guest;
\end{verbatim}
\normalsize

\vspace{0.1cm}
\footnotesize
DCL — практическая реализация функции \textbf{безопасность}:
студент из примера про безопасность (раздел о функциях СУБД)
видит свои оценки именно благодаря правам, выданным через DCL.

\end{frame}


\begin{frame}[fragile]{TCL — язык управления транзакциями}

\textbf{TCL (Transaction Control Language)} — команды для
управления границами транзакции.

\vspace{0.3cm}

\small
\begin{table}
\centering
\begin{tabular}{>{\ttfamily}p{0.22\textwidth}|p{0.62\textwidth}}
\toprule
\textnormal{Команда} & \textnormal{Назначение}\\
\midrule
BEGIN & начать транзакцию\\
COMMIT & зафиксировать все изменения транзакции\\
ROLLBACK & отменить все изменения транзакции\\
\bottomrule
\end{tabular}
\end{table}

\vspace{0.3cm}

\begin{verbatim}
BEGIN;
UPDATE student SET group_id = 2 WHERE id = 5;
UPDATE student_group SET name = 'ПМИ-25' WHERE id = 2;
COMMIT;
\end{verbatim}
\normalsize

\vspace{0.1cm}
\footnotesize
Обе команды \texttt{UPDATE} выполнятся как единое целое
(«всё или ничего») — это и есть \textbf{атомарность}
транзакции, упомянутая при определении термина.

\end{frame}


\begin{frame}{Предпросмотр: что дальше в лабораторных работах}

Сегодня мы увидели базовый синтаксис DDL, DML, DCL, TCL
для работы с одной таблицей.

\vspace{0.3cm}

На лабораторных работах и в следующих лекциях появятся:

\begin{itemize}
  \item выборка данных сразу из нескольких таблиц — \texttt{JOIN};
  \item агрегатные функции — \texttt{COUNT}, \texttt{SUM},
        \texttt{AVG}, \texttt{GROUP BY};
  \item вложенные запросы (подзапросы);
  \item индексы для ускорения запросов (лекция о физическом
        проектировании).
\end{itemize}

\vspace{0.3cm}

\begin{alertblock}{Важно}
Синтаксис, показанный сегодня, — рабочий минимум для первой
лабораторной работы; более сложные запросы разбираются
отдельно.
\end{alertblock}

\end{frame}


% =========================================================
% РАЗДЕЛ 8. ФИЗИЧЕСКОЕ ХРАНЕНИЕ ДАННЫХ
% =========================================================

\begin{frame}{От логической модели к физическому хранению}

\begin{center}
\resizebox{0.98\textwidth}{!}{%
\begin{tikzpicture}[node distance=1cm]

\node[entity, minimum width=3.1cm] (log) {Логическая структура\\(таблицы, столбцы)};
\node[catbox, right=of log] (files) {Файлы};
\node[catbox, right=of files] (pages) {Страницы};
\node[catbox, right=of pages, minimum width=2.6cm] (buf) {Буфер\\(кэш в ОЗУ)};

\draw[-{Latex[length=2.5mm]}, thick, svfuAccent] (log) -- (files);
\draw[-{Latex[length=2.5mm]}, thick, svfuAccent] (files) -- (pages);
\draw[-{Latex[length=2.5mm]}, thick, svfuAccent] (pages) -- (buf);

\end{tikzpicture}%
}
\end{center}

\footnotesize
Пользователь и приложение работают только с логической
структурой; всё, что ниже, скрыто внутренним уровнем
архитектуры ANSI/SPARC.

\end{frame}


\begin{frame}{Файлы, страницы, буферы — кратко}

\begin{itemize}
  \item \textbf{Файл} — единица хранения таблицы или её части
        на диске;
  \item \textbf{Страница} — блок фиксированного размера,
        которым СУБД читает и записывает данные на диск;
  \item \textbf{Буфер} — область оперативной памяти,
        где хранятся недавно использованные страницы,
        чтобы не обращаться к диску повторно.
\end{itemize}

\vspace{0.3cm}

\begin{alertblock}{Важно}
Это обзорное представление. Алгоритмы работы буфера
и физической организации файлов — предмет отдельной
будущей лекции о физическом проектировании.
\end{alertblock}

\end{frame}


\begin{frame}{Индекс как идея}

Поиск студента по фамилии в таблице из миллиона строк:

\vspace{0.3cm}

\begin{itemize}
  \item \textbf{без индекса} — СУБД просматривает
        все строки подряд;
  \item \textbf{с индексом} — СУБД использует
        вспомогательную структуру, которая сразу указывает,
        где искать нужные строки.
\end{itemize}

\vspace{0.4cm}

\begin{block}{Определение}
\textbf{Индекс} — вспомогательная структура данных,
ускоряющая поиск записей ценой дополнительного
объёма хранения и дополнительных затрат при изменении
данных.
\end{block}

\vspace{0.2cm}
\footnotesize
Синтаксис создания индекса и стратегии индексирования —
предмет отдельной будущей лекции.

\end{frame}


% =========================================================
% РАЗДЕЛ 9. ИТОГИ
% =========================================================

\begin{frame}{Итоговая схема: от файла к системе базы данных}

\begin{center}
\resizebox{0.98\textwidth}{!}{%
\begin{tikzpicture}[node distance=0.5cm and 0.4cm]

\node[phase] (file) {Файл};
\node[phase, right=of file] (table) {Таблица};
\node[phase, right=of table] (db) {База\\данных};
\node[phase, right=of db] (dbms) {СУБД};
\node[phase, right=of dbms] (sys) {Система\\базы данных};

\draw[-{Latex[length=2.5mm]}, thick, svfuAccent] (file) -- (table);
\draw[-{Latex[length=2.5mm]}, thick, svfuAccent] (table) -- (db);
\draw[-{Latex[length=2.5mm]}, thick, svfuAccent] (db) -- (dbms);
\draw[-{Latex[length=2.5mm]}, thick, svfuAccent] (dbms) -- (sys);

\end{tikzpicture}%
}
\end{center}

\vspace{0.4cm}

\begin{alertblock}{Главная идея лекции}
СУБД — не просто программа для хранения таблиц,
а комплекс функций (транзакции, целостность,
многопользовательский доступ, восстановление,
безопасность), реализованных набором программных
компонентов и доступных через SQL.
\end{alertblock}

\end{frame}


\begin{frame}{Что будет раскрыто подробнее в следующих лекциях}

\small
\begin{table}
\centering
\begin{tabular}{p{0.5\textwidth}|p{0.4\textwidth}}
\toprule
\textbf{Тема} & \textbf{Где}\\
\midrule
Реляционная модель, ключи, ограничения целостности (SQL) & Лекция 3\\
Нормализация и нормальные формы & Лекция 3 / отдельная лекция\\
Сложные запросы: \texttt{JOIN}, агрегатные функции, подзапросы & SQL-практикум\\
Механика транзакций: изоляция, блокировки & Отдельная лекция о транзакциях\\
Детали журналирования и восстановления & Та же лекция о надёжности\\
Синтаксис и стратегии индексирования & Лекция о физическом проектировании\\
\bottomrule
\end{tabular}
\end{table}

\end{frame}


\begin{frame}{Контрольные вопросы}

\small
\begin{enumerate}
  \item В чём разница между базой данных, СУБД и системой базы данных?
  \item Что различают схема и экземпляр базы данных?
  \item Чем архитектура ANSI/SPARC отличается от внутренней архитектуры СУБД?
  \item Назовите пять программных компонентов ядра СУБД.
  \item Что такое клиент-серверная модель работы с СУБД?
  \item Чем отличаются DDL, DML, DCL и TCL — приведите по одной команде на каждый.
  \item Какая команда SQL отвечает за атомарное выполнение нескольких операций?
  \item Зачем нужен индекс и какова его цена?
\end{enumerate}

\end{frame}


\begin{frame}[fragile]{Задание к следующей лекции}

Возьмите предметную область и ER-диаграмму,
разработанную к Лекции 1.

\vspace{0.3cm}

Напишите:
\begin{enumerate}
  \item команды \texttt{CREATE TABLE} (DDL) для 2--3 сущностей
        вашей предметной области;
  \item по одной команде \texttt{INSERT}, \texttt{SELECT},
        \texttt{UPDATE} (DML) на эти таблицы.
\end{enumerate}

\vspace{0.3cm}

\textbf{Результат:} черновой набор DDL/DML-команд —
основа для первой лабораторной работы.

\end{frame}


% --- Спасибо ---
\begin{frame}[plain,noframenumbering]

\thispagestyle{empty}

\vspace*{\fill}

\begin{center}

{\LARGE\bfseries
Спасибо за внимание!
}

\end{center}

\vspace*{\fill}

\end{frame}


\end{document}